% !TeX root = ../thuthesis-example.tex

\begin{abstract}
  本文围绕电商用户行为数据分析课题，构建了一套基于 Docker ARM64、Hadoop HDFS/YARN、Spark SQL、Flask API 与 React 前端的本机大数据分析实验系统。系统以 Kaggle 电商行为数据为主要输入，完成原始 CSV 数据采样、HDFS 上传、Spark 分布式读取、字段标准化、异常值过滤、重复事件去重、会话事实表生成、指标缓存发布和前端可视化展示。业务层面实现了基础看板、转化漏斗、用户行为路径（含事件转移矩阵）、特征集市、需求预测（双尺度 AR-MLP）、商品关联图谱（Support/Confidence/Lift）、Cohort 留存、商品组合分析（HHI 浓缩度）、购物车挽回、收入归因、运营优化（MILP 与贪心降级）、个性化推荐（含 LR/GBT 排序器与 precision/recall/NDCG 评估指标）、异常检测（Robust Z-Score）、生命周期分层、受控自然语言查询、实验分流（含 SRM 卡方检验与 AUUC 评估）、决策引擎（决策自动生成）和 Ops 证据面板等模块。各模块统一复用清洗后的事件表、会话事实表和日粒度特征表，保证不同页面中的行数、收入、转化率、留存率、推荐置信度和异常分级口径一致。

  在工程实现上，本文重点解决了本机 ARM Mac 环境下 Hadoop/Spark 集群复现、端口冲突、跨架构镜像不可用、Spark History 日志过大、driver 端全量 collect 造成内存风险等问题。系统采用 YARN client mode 作为第一阶段提交方式，保留 Flask 读取本地缓存的便利性；同时使用 AQE、合理 shuffle 分区、\texttt{MEMORY\_AND\_DISK} 缓存、Parquet 清洗层、preview JSON 上界和模块级质量门禁降低运行风险。实验部分对比 local Spark、YARN-only、YARN+AQE、YARN+算法限流和 YARN+Parquet 多组配置，并补充 1\%/5\% 用户级样本与典型模块 benchmark。结果表明，仅把任务提交到 YARN 并不能自动提升性能或避免 OOM；真正有效的优化来自 Spark SQL 执行计划、列式数据布局、算法输出边界和可追溯的运行证据。最终系统能够在课程设计环境中完成从分布式存储、Spark 计算、后端接口到前端可视化和报告分析的完整闭环。
  \thusetup{
    keywords = {Hadoop YARN, Spark SQL, Docker ARM64, 电商行为分析, 内存优化, Benchmark},
  }
\end{abstract}
